% !PRINTED 22
\weberpage{22}
Daraus aber erkennt man leicht, dass diese Gruppen nicht nur von
gleichem Grade sind, sondern dass sie \gesperrt{isomorph} sind. Denn es ist
gleichzeitig
\[
Rb_1\,Rb_2 = Rb_1b_2
\]
und
\[
Q'b_1\,Q'b_2 = Q'b_1b_2.
\]
Ebenso kann man auch schliessen, dass die beiden Gruppen
\[
Q'/R,\quad P/Q
\]
isomorph sind.

Die Gruppen $P/Q$ und $P/Q'$ sind aber, da $Q$, $Q'$ grösste
Normaltheiler sind, nach 3) einfach, und folglich sind auch die
damit isomorphen Gruppen $Q'/R$ und $Q/R$ einfach, und folglich
haben wir den Satz:

\gesperrt{4. Sind $Q$ und $Q'$ zwei verschiedene grösste
Normaltheiler von $P$ und ist $R$ der Durchschnitt von $Q$ und
$Q'$, so ist $R$ grösster Normaltheiler sowohl von $Q$ als von
$Q'$, und der Index von $R$ in Bezug auf $Q$ ist gleich dem Index
von $Q'$ in Bezug auf $P$.}

Damit haben wir die Grundlage gewonnen, um sehr einfach den Satz
von der Constanz der Indexreihe nachzuweisen.

Wir schreiben zur besseren Uebersicht die verschiedenen in
Vergleich zu ziehenden Compositionsreihen so, dass wir den Index
eines jeden Gliedes in Bezug auf das vorangehende unter das Glied
setzen.

Danach mögen irgend zwei gegebene Compositionsreihen von $P$ mit
den zugehörigen Indices folgende sein:
\weq{23}{\begin{aligned}
&P,\ Q,\ Q_1,\ Q_2.\ Q_3 \ldots \\ % [sic: the source prints a full stop after Q_2.]
&\hphantom{P,\ }\nu,\ \nu_1,\ \nu_2,\ \nu_3 \ldots
\end{aligned}}
\weq{24}{\begin{aligned}
&P,\ Q',\ Q_1',\ Q_2',\ Q_3' \ldots \\
&\hphantom{P,\ }\nu',\ \nu_1',\ \nu_2',\ \nu_3' \ldots
\end{aligned}}
Es sei nun $R$ der Durchschnitt von $Q$ und $Q'$, und $\mu$ und
$\mu'$ seien die Indices von $R$ in Bezug auf $Q$ und $Q'$. Wegen
des Isomorphismus der Gruppen $P/Q$ und $Q'/R$ ist dann $\mu' = \nu$,
und aus dem entsprechenden Grunde $\mu = \nu'$. Wir bilden eine
Compositionsreihe von $R$ mit der zugehörigen Indexreihe
\weq{25}{\begin{aligned}
&R,\ R_1,\ R_2,\ R_3 \ldots \\
&\hphantom{R,\ }\mu_1,\ \mu_2,\ \mu_3 \ldots,
\end{aligned}}
%CONT>>
